\section{Main results}\label{sec:results}

Each statement below is labelled with the Lean declaration that proves it. ``Formal'' means machine-checked,
with axioms \lean{propext}, \lean{Classical.choice} and \lean{Quot.sound} only, and with
\eqref{eq:grh} discharged by the GRH proof (Section~\ref{sec:formal}). Statements marked ``informal'' are
proved in this paper but are not formalized.

\subsection{A uniform bound via Borel--Carath\'eodory}

\begin{theorem}[formal; \lean{GRHPseudo.Main.primeCharSum}]\label{thm:A}
For every $q\ge1$, every $\chi\ne\chi_0$ modulo $q$ (primitive or not) and every integer $N\ge3$,
\[
  \Bigl\lvert\sum_{0<n<N}\chi(n)\Lambda(n)\Bigl(1-\frac nN\Bigr)\Bigr\rvert\le\frac{48e}{\pi}(\log N)^2\sqrt N\,
  \bigl(\pi\log(2+8q)+J\bigr),\qquad J=\int_\R\frac{\log(4+\abs t)}{1+t^2}\,dt .
\]
\end{theorem}

\begin{corollary}[formal; \lean{GRHPseudo.Main.bias_primeMeasure}]\label{cor:A}
Let $S$ be the set of $n\in(0,N)$ coprime to $q$, with weights $\Lambda(n)(1-n/N)$. For $N\ge3$ with
$L:=\frac12\bigl(\lfloor N/2\rfloor\log2-\log(\lfloor N/2\rfloor+1)\bigr)-\log q\,\lfloor\log_2N\rfloor>0$, the
bias of $S$ against every $\chi\ne\chi_0$ is at most the right side of Theorem~\ref{thm:A} divided by $L$.
\end{corollary}

The constants here are poor: to certify bias $\le1/2$ for $q\le2^{64}$ one needs $N\approx9.7\cdot10^{15}$
(Table~\ref{tab:params}). The loss is intrinsic to the route. Borel--Carath\'eodory bounds $L'/L$
pointwise on $\operatorname{Re}s=\frac12+1/\log N$ only by $O((\log N)^2\log q)$, whereas the explicit formula
avoids pointwise bounds altogether.

\subsection{The sharp bound}

Fix $a\ge0$ and $y>2a$, put $x=e^y$ and $M=\lfloor e^{y+2a}\rfloor+1$, and define the sample space
\[
  S_q(a,y)=\{0\le n<M:\gcd(n,q)=1\},\qquad w_{a,y}(n)=\frac{\Lambda(n)\,k_a(\log n-y)}{\sqrt n}.
\]
Its weights are nonnegative and supported on prime powers in $[xe^{-2a},xe^{2a}]$. With $f=f_a$, put
\begin{align*}
  I_a(\beta)&=\int_\R\abs{\hat f_a(r)}^2\bigl(\log(4+\abs r)+\beta\bigr)\,dr,\qquad
  \mathrm{SP}(a)=\sum_{n\le e^{2a}}\frac{\Lambda(n)}{\sqrt n}\bigl(\abs{k_a(\log n)}+\abs{k_a(-\log n)}\bigr),\\
  U_q(a,y)&=\sum_{n<M,\ \gcd(n,q)>1}\abs{w_{a,y}(n)} ,
\end{align*}
and
\begin{align*}
  \mathrm{Num}&=k_a(0)\log\tfrac q\pi+\mathrm{SP}(a)+\frac{I_a(5)+I_a(\frac12)}{2\pi}+U_q(a,y),\\
  \mathrm{Den}&=e^{y/2}A(a)-B(a)-\frac{I_a(5+\log\pi)+I_a(\frac12)}{2\pi}-U_q(a,y).
\end{align*}

\begin{theorem}[formal; \lean{GRHPseudo.Sharp.bias_sharp}, \lean{sharp_charSum_mass}]\label{thm:sharp}
For every $q\ge1$, every $\chi\ne\chi_0$ modulo $q$, every $a\ge0$ and every $y>2a$:
\[
  \Bigl\lvert\sum_{n\in S_q(a,y)}w_{a,y}(n)\chi(n)\Bigr\rvert\le\mathrm{Num},\qquad
  \sum_{n\in S_q(a,y)}w_{a,y}(n)\ge\mathrm{Den}.
\]
In particular, if $\mathrm{Den}>0$ then $\bias_\chi(S_q(a,y),w_{a,y})\le\mathrm{Num}/\mathrm{Den}$.
\end{theorem}

The leading term of $\mathrm{Num}/\mathrm{Den}$ is $k_a(0)\log q/(\sqrt x\,A(a))$. All remaining terms are
either independent of $q$ or $O(\log q\log x/\sqrt x)$.

\begin{remark}
This leading term agrees with the asymptotic Lemma~6.1 of Lamzouri--Li--Soundararajan \cite{lls2015} (see
Section~\ref{sec:related}). Theorem~\ref{thm:sharp} adds the following:
\begin{itemize}
\item every lower-order term is explicit, with no $o(1)$;
\item the numerator bounds $\abs{\sum w\chi}$ itself, not just its real part, and holds for all $\chi\ne\chi_0$, primitive or not;
\item the mass bound is explicit, obtained from the explicit formula for $\zeta$ under RH;
\item the whole statement is machine-checked.
\end{itemize}
The explicit form makes small certificates possible: with an $o(1)$ one cannot conclude anything at any specific $q$.
\end{remark}

\subsection{Closed form}
Let $K(a)=\sinh a+a\cosh a$, $C(a)=(4\sinh\frac a2)^2+(2\sinh\frac a2+2\cosh\frac a2)^2$,
$U_B=K(a)e^{a-y/2}\log q\,\lfloor\log_2M\rfloor$, and
\begin{align*}
  \mathrm{num}_B&=(\sinh a+a)\log q+\frac{4K(a)}{e}\bigl(\lfloor e^{2a}\rfloor+1\bigr)
     +\frac{C(a)\bigl(12+\frac{11}2\pi\bigr)}{2\pi}+U_B,\\
  \mathrm{den}_B&=e^{y/2}\bigl(\sinh a'+a'\bigr)^2-2(\sinh a+a)^2-\frac{C(a)\bigl(12+(\frac{11}2+\log\pi)\pi\bigr)}{2\pi}-U_B .
\end{align*}

\begin{theorem}[formal; \lean{GRHPseudo.Main.bias_sharp_closed}]\label{thm:closed}
For every $q\ge1$, $\chi\ne\chi_0$ modulo $q$, $a\ge1$ and $y>2a$ with $\mathrm{den}_B>0$,
\[
  \bias_\chi\bigl(S_q(a,y),w_{a,y}\bigr)\le\frac{\mathrm{num}_B}{\mathrm{den}_B}.
\]
\end{theorem}

\begin{corollary}[formal; \lean{GRHPseudo.Main.bias_le_half_of_le_two_pow_64}]\label{cor:cert}
Let $a=1$ and $y=2\log150$, so $x=22500$ and $S_q(1,y)$ consists of the integers below
$\lfloor22500e^2\rfloor+1=166254$ coprime to $q$. For every $1\le q\le2^{64}$ and every nontrivial
Dirichlet character $\chi$ modulo $q$,
\[
  \bias_\chi\bigl(S_q(1,y),w_{1,y}\bigr)\le\tfrac12 .
\]
Equivalently, every nontrivial eigenvalue of the weighted Cayley graph $\mathrm{Cay}(\Gq,w_{1,y})$ has modulus
at most $1/2$. (The proof gives $\mathrm{num}_B\le236.6$ and $\mathrm{den}_B\ge555.8$, hence bias $\le0.4257$.)
\end{corollary}

\subsection{The leading constant}

\begin{proposition}[informal; abstract core formal: \lean{GRHPseudo.inner_mul_inner_le}]\label{prop:opt}
Let $a>0$ and let $f\in L^2[-a,a]$ be real and nonzero, with $k=f*\tilde f$. Then
\[
  \frac{k(0)}{\int_\R k(u)e^{u/2}\,du}\ \ge\ \frac1{a+\sinh a},
\]
with equality if and only if $f$ is a nonzero multiple of $\cosh(t/2)\mathbf 1_{[-a,a]}$. By the
Boas--Kac--Krein factorization \cite{boaskac1945,ehmgneitingrichards2004}, every continuous positive-definite
$k$ supported in $[-2a,2a]$ is of this form.
\end{proposition}

The quantity in Proposition~\ref{prop:opt} is the coefficient of $\log q/\sqrt x$ that the positivity method
produces from a \emph{positive-definite} kernel supported in $[-2a,2a]$ (Section~\ref{sec:optimality}). It is not
the limit of the method: non-even kernels whose Fourier transform is dominated by a positive-definite one do better.
Normalizing by the top $X$ of the support:

\begin{corollary}[informal]\label{cor:asymp}
Write the leading term of the positivity bound, or of the bound of \cite[Lemma~6.1]{lls2015}, as $c\log q/\sqrt X$, where
$X$ is the top of the support.
\begin{enumerate}
\item \emph{(Barrier.)} Every kernel gives $c\ge1$ (Proposition~\ref{prop:barrier}).
\item \emph{(Positive-definite kernels.)} A positive-definite kernel supported in an interval of length $4a$ gives
  $c\ge e^a/(a+\sinh a)\ge c^*=2/a_*^2=1.22363\ldots$, where $a_*=1.27846\ldots$ is the root of $a=1+e^{-a}$.
\item \emph{(Achievability.)} For every $\delta>0$ there is a smooth nonnegative convolution-square kernel with
  $\bias\le(1+\delta+o(1))\log q/\sqrt X$ as $q\to\infty$ and $X/\log^2q\to\infty$. Hence support
  $X\ge(1+o(1))\log^2q/\varepsilon^2$ suffices for bias $\le\varepsilon$, and within the method nothing smaller does.
\end{enumerate}
For the formally verified kernel $f_a$ (positive-definite, transition width $1/64$), Theorem~\ref{thm:closed} gives leading
coefficient at most $e^a(\sinh a+a)/(\sinh a'+a')^2$, which equals $1.2626$ at $a=a_*$ and $1.2964$ at $a=1$.
\end{corollary}

\subsection{An explicit bound with a convolution-square kernel}
Corollary~\ref{cor:asymp} says that a non-even kernel beats every positive-definite one. We have made this explicit
and formal for one family. Let $A>0$, let $S$ be Mathlib's \lean{Real.smoothTransition}, and put
\[
  f_A(t)=e^{t/2}\,S\bigl(64(t+A)\bigr)\,S(-64t),\qquad K_A=f_A*f_A,\qquad M_A=f_A*\tilde f_A .
\]
So $f_A\ge0$ is supported in $[-A,0]$ and equals $e^{t/2}$ on $[-A+\frac1{64},-\frac1{64}]$, $K_A\ge0$ is supported in
$[-2A,0]$, and $\abs{\hat K_A}=\abs{\hat f_A}^2=\hat M_A$ on $\R$ (Proposition~\ref{prop:convsq}). For $y>2A$ and $x=e^y$
let $S^\square_q(y)=\{0\le n\le\lfloor x\rfloor:\gcd(n,q)=1\}$ with weights
$w^\square_{A,y}(n)=\Lambda(n)K_A(\log n-y)n^{-1/2}$, which are supported on prime powers in $[xe^{-2A},x]$. For $R>0$
put $G(B)=(\log(4+R)+B)(1-e^{-A})+36/(\pi R)$ and $U^\square_B=Ae^{-y/2}\log q\,\lfloor\log_2(\lfloor x\rfloor+1)\rfloor$, and
\begin{align*}
  \mathrm{num}^\square_B&=(1-e^{-A})\log q+\tfrac2e\bigl(\lfloor e^{A}\rfloor+1\bigr)+G(5)+G(\tfrac12)+U^\square_B,\\
  \mathrm{den}^\square_B&=e^{y/2}\bigl(e^{-1/64}-e^{-A+1/64}\bigr)^2-2A-G(5+\log\pi)-G(\tfrac12)-U^\square_B .
\end{align*}

\begin{theorem}[formal; \lean{GRHPseudo.Main.bias_convSq_closed}]\label{thm:convclosed}
For every $q\ge1$, $\chi\ne\chi_0$ modulo $q$, $A\ge\frac1{32}$, $R>0$ and $y>2A$ with $\mathrm{den}^\square_B>0$,
\[
  \bias_\chi\bigl(S^\square_q(y),w^\square_{A,y}\bigr)\le\frac{\mathrm{num}^\square_B}{\mathrm{den}^\square_B}.
\]
\end{theorem}

The proof runs the argument of Theorem~\ref{thm:sharp} with the zero sum of the twisted $K_A$ bounded by that of $M_A$
(\lean{norm_primeSide_twist_le_maj}, \lean{primeSideZeta_twist_eq_maj}); the intermediate bound with exact constants is
\lean{GRHPseudo.Main.bias_convSq_sharp}. It also uses a sharper gamma-side estimate (\lean{Ilog_div_le_sharp}): if
$f\in C^2_c$ and $\abs r\,\abs{\hat f(r)}\le C$, then for $B\ge0$ and $R>0$
\[
  \frac1{2\pi}\int_\R\abs{\hat f(r)}^2\bigl(\log(4+\abs r)+B\bigr)dr\le\bigl(\log(4+R)+B\bigr)(f*\tilde f)(0)+\frac{4C^2}{\pi R}.
\]
Here $C=3$, from $\int\abs{f_A'}\le3$. The leading coefficient of $\mathrm{num}^\square_B/\mathrm{den}^\square_B$, normalized by the top
$x$ of the support, is $(1-e^{-A})/(e^{-1/64}-e^{-A+1/64})^2$, which equals $1.162$ at $A=\frac94$; for the formal $\cosh$ kernel
the corresponding value at $a=1$ is $1.2964$.

\begin{corollary}[formal; \lean{GRHPseudo.Main.bias_le_half_of_le_two_pow_64_convSq}]\label{cor:convcert}
Let $A=\frac94$, $R=12$ and $y=2\log220$, so $x=48400$ and $S^\square_q(y)$ consists of the integers $n\le48400$ coprime to $q$.
For every $1\le q\le2^{64}$ and every nontrivial Dirichlet character $\chi$ modulo $q$,
\[
  \bias_\chi\bigl(S^\square_q(y),w^\square_{9/4,y}\bigr)\le\tfrac12 .
\]
(The proof gives $\mathrm{num}^\square_B\le65.7$ and $\mathrm{den}^\square_B\ge143.9$, hence bias $\le0.4566$.) The support is $3.4$ times
smaller than in Corollary~\ref{cor:cert}.
\end{corollary}
